\hypertarget{os-services-signal-handling-and-subprocesses}{%
\chapter{OS Services, Signal Handling, and
Subprocesses}\label{os-services-signal-handling-and-subprocesses}}

Python is widely used for system administration and process
orchestration because it provides a near-one-to-one mapping to operating
system primitives while managing the complexity of cross-platform
differences.

\hypertarget{os-the-system-call-bridge}{%
\subsection{\texorpdfstring{37.1 \texttt{os}: The System Call
Bridge}{37.1 os: The System Call Bridge}}\label{os-the-system-call-bridge}}

The \passthrough{\lstinline!os!} module is the primary interface to the
OS. * \textbf{System Calls}: Most functions in
\passthrough{\lstinline!os!} are thin wrappers around C standard library
calls (\passthrough{\lstinline!open!}, \passthrough{\lstinline!read!},
\passthrough{\lstinline!write!}, \passthrough{\lstinline!fork!},
\passthrough{\lstinline!exec!}). * \textbf{Environment}:
\passthrough{\lstinline!os.environ!} is a mapping object that syncs with
the process's environment block. * \textbf{Filesystem}: Provides
low-level manipulation of file descriptors
(\passthrough{\lstinline!os.dup!}, \passthrough{\lstinline!os.pipe!})
and path metadata (\passthrough{\lstinline!os.stat!}).

\hypertarget{io-the-stream-hierarchy}{%
\subsection{\texorpdfstring{37.2 \texttt{io}: The Stream
Hierarchy}{37.2 io: The Stream Hierarchy}}\label{io-the-stream-hierarchy}}

The \passthrough{\lstinline!io!} module provides the foundations for
Python's file handling. It uses a tiered architecture: 1. \textbf{Raw
I/O}: \passthrough{\lstinline!FileIO!} objects represent raw OS file
descriptors. They perform unbuffered system calls. 2. \textbf{Buffered
I/O}: \passthrough{\lstinline!BufferedReader!},
\passthrough{\lstinline!BufferedWriter!}, and
\passthrough{\lstinline!BufferedRandom!}. These maintain internal
C-level buffers to minimize the number of expensive system calls. 3.
\textbf{Text I/O}: \passthrough{\lstinline!TextIOWrapper!} handles
encoding and decoding (e.g., UTF-8 to Unicode) on top of a buffered
binary stream.

\hypertarget{signal-asynchronous-event-handling}{%
\subsection{\texorpdfstring{37.3 \texttt{signal}: Asynchronous Event
Handling}{37.3 signal: Asynchronous Event Handling}}\label{signal-asynchronous-event-handling}}

Signals are software interrupts sent to a process. Handling them in a
virtual machine like CPython is complex.

\hypertarget{the-main-thread-restriction}{%
\subsubsection{1. The Main Thread
Restriction}\label{the-main-thread-restriction}}

In Python, signals are always received by the main thread, regardless of
which thread was executing when the signal arrived. *
\textbf{Internals}: When a signal arrives, the OS interrupts the
process. The C-level signal handler in CPython sets a flag. *
\textbf{The Check}: The evaluation loop
(\passthrough{\lstinline!ceval.c!}) periodically checks this flag. If
set, it invokes the Python-level signal handler. This ensures that
Python code only runs at ``safe'' points where the VM state is
consistent.

\hypertarget{signal.set_wakeup_fd}{%
\subsubsection{\texorpdfstring{2.
\texttt{signal.set\_wakeup\_fd}}{2. signal.set\_wakeup\_fd}}\label{signal.set_wakeup_fd}}

For integration with event loops (like
\passthrough{\lstinline!asyncio!}),
\passthrough{\lstinline!set\_wakeup\_fd!} writes a byte to a file
descriptor whenever a signal is received, allowing a
\passthrough{\lstinline!select()!} or \passthrough{\lstinline!poll()!}
call to wake up and handle the signal.

\hypertarget{subprocess-orchestrating-external-processes}{%
\subsection{\texorpdfstring{37.4 \texttt{subprocess}: Orchestrating
External
Processes}{37.4 subprocess: Orchestrating External Processes}}\label{subprocess-orchestrating-external-processes}}

The \passthrough{\lstinline!subprocess!} module is the modern
replacement for \passthrough{\lstinline!os.system!} and
\passthrough{\lstinline!os.spawn!}.

\hypertarget{process-creation}{%
\subsubsection{1. Process Creation}\label{process-creation}}

\begin{itemize}
\tightlist
\item
  \textbf{POSIX}: Uses \passthrough{\lstinline!fork()!} and
  \passthrough{\lstinline!exec()!}. Between the fork and exec, it
  handles closing file descriptors and setting up pipes.
\item
  \textbf{Windows}: Uses the \passthrough{\lstinline!CreateProcess()!}
  API.
\end{itemize}

\hypertarget{pipe-multiplexing}{%
\subsubsection{2. Pipe Multiplexing}\label{pipe-multiplexing}}

\passthrough{\lstinline!subprocess.communicate()!} reads data from
\passthrough{\lstinline!stdout!} and \passthrough{\lstinline!stderr!}
while writing to \passthrough{\lstinline!stdin!}. * \textbf{Deadlock
Prevention}: It uses \passthrough{\lstinline!selectors!} (or threads on
Windows) to read from multiple pipes simultaneously. This prevents the
``buffer full'' deadlock that occurs if you try to read from one pipe
while the process is blocked trying to write to another.

\hypertarget{shlex-safe-command-parsing}{%
\subsubsection{\texorpdfstring{3. \texttt{shlex}: Safe Command
Parsing}{3. shlex: Safe Command Parsing}}\label{shlex-safe-command-parsing}}

When building command strings, always use
\passthrough{\lstinline!shlex.split()!} to ensure that arguments with
spaces or special characters are handled correctly, preventing shell
injection vulnerabilities.

\begin{center}\rule{0.5\linewidth}{0.5pt}\end{center}

\begin{center}\rule{0.5\linewidth}{0.5pt}\end{center}
